\section{The spectrum of a ring}
\label{section-spectrum-ring}

\noindent
We arbitrarily decide that the spectrum of a ring as a topological space
is part of the algebra chapter, whereas an affine scheme is part of the
chapter on schemes.

\begin{definition}
\label{definition-spectrum-ring}
Let $R$ be a ring.
\begin{enumerate}
\item The {\it spectrum} of $R$ is the set of prime ideals of $R$.
It is usually denoted $\Spec(R)$.
\item Given a subset $T \subset R$ we let $V(T) \subset \Spec(R)$
be the set of primes containing $T$, i.e., $V(T) = \{ \mathfrak p \in
\Spec(R) \mid \forall f\in T, f\in \mathfrak p\}$.
\item Given an element $f \in R$ we let $D(f) \subset \Spec(R)$
be the set of primes not containing $f$.
\end{enumerate}
\end{definition}

\begin{lemma}
\label{lemma-Zariski-topology}
Let $R$ be a ring.
\begin{enumerate}
\item The spectrum of a ring $R$ is empty if and only if $R$
is the zero ring.
\item Every nonzero ring has a maximal ideal.
\item Every nonzero ring has a minimal prime ideal.
\item Given an ideal $I \subset R$ and a prime ideal
$I \subset \mathfrak p$ there exists a prime
$I \subset \mathfrak q \subset \mathfrak p$ such
that $\mathfrak q$ is minimal over $I$.
\item If $T \subset R$, and if $(T)$ is the ideal generated by
$T$ in $R$, then $V((T)) = V(T)$.
\item If $I$ is an ideal and $\sqrt{I}$ is its radical,
see basic notion (\ref{item-radical-ideal}), then $V(I) = V(\sqrt{I})$.
\item Given an ideal $I$ of $R$ we have $\sqrt{I} =
\bigcap_{I \subset \mathfrak p} \mathfrak p$.
\item If $I$ is an ideal then $V(I) = \emptyset$ if and only
if $I$ is the unit ideal.
\item If $I$, $J$ are ideals of $R$ then $V(I) \cup V(J) =
V(I \cap J)$.
\item If $(I_a)_{a\in A}$ is a set of ideals of $R$ then
$\bigcap_{a\in A} V(I_a) = V(\bigcup_{a\in A} I_a)$.
\item If $f \in R$, then $D(f) \amalg V(f) = \Spec(R)$.
\item If $f \in R$ then $D(f) = \emptyset$ if and only if $f$
is nilpotent.
\item If $f = u f'$ for some unit $u \in R$, then $D(f) = D(f')$.
\item If $I \subset R$ is an ideal, and $\mathfrak p$ is a prime of
$R$ with $\mathfrak p \not\in V(I)$, then there exists an $f \in R$
such that $\mathfrak p \in D(f)$, and $D(f) \cap V(I) = \emptyset$.
\item If $f, g \in R$, then $D(fg) = D(f) \cap D(g)$.
\item If $f_i \in R$ for $i \in I$, then
$\bigcup_{i\in I} D(f_i)$ is the complement of $V(\{f_i \}_{i\in I})$
in $\Spec(R)$.
\item If $f \in R$ and $D(f) = \Spec(R)$, then $f$ is a unit.
\end{enumerate}
\end{lemma}

\begin{proof}
We address each part in the corresponding item below.
\begin{enumerate}
\item This is a direct consequence of (2) or (3).
\item Let $\mathfrak{A}$ be the set of all proper ideals of $R$. This set is
ordered by inclusion and is non-empty, since $(0) \in \mathfrak{A}$ is a proper
ideal. Let $A$ be a totally ordered subset of $\mathfrak A$.
Then $\bigcup_{I \in A} I$ is in
fact an ideal. Since 1 $\notin I$ for all $I \in A$, the union does not contain
1 and thus is proper. Hence $\bigcup_{I \in A} I$ is in $\mathfrak{A}$ and is
an upper bound for the set $A$. Thus by Zorn's lemma $\mathfrak{A}$ has a
maximal element, which is the sought-after maximal ideal.
\item Since $R$ is nonzero, it contains a maximal ideal which is a prime ideal.
Thus the set $\mathfrak{A}$ of all prime ideals of $R$ is nonempty.
$\mathfrak{A}$ is ordered by reverse-inclusion. Let $A$ be a totally ordered
subset of $\mathfrak{A}$. It's pretty clear that $J = \bigcap_{I \in A} I$ is
in fact an ideal. Not so clear, however, is that it is prime. Let $xy \in J$.
Then $xy \in I$ for all $I \in A$. Now let $B = \{I \in A | y \in I\}$. Let $K
= \bigcap_{I \in B} I$. Since $A$ is totally ordered, either $K = J$ (and we're
done, since then $y \in J$) or $K \supset J$ and for all $I \in A$ such that
$I$ is properly contained in $K$, we have $y \notin I$. But that means that for
all those $I, x \in I$, since they are prime. Hence $x \in J$. In either case,
$J$ is prime as desired. Hence by Zorn's lemma we get a maximal element which
in this case is a minimal prime ideal.
\item This is the same exact argument as (3) except you only consider prime
ideals contained in $\mathfrak{p}$ and containing $I$.
\item $(T)$ is the smallest ideal containing $T$. Hence if $T \subset I$, some
ideal, then $(T) \subset I$ as well. Hence if $I \in V(T)$, then $I \in V((T))$
as well. The other inclusion is obvious.
\item Since $I \subset \sqrt{I}, V(\sqrt{I}) \subset V(I)$. Now let
$\mathfrak{p} \in V(I)$. Let $x \in \sqrt{I}$. Then $x^n \in I$ for some $n$.
Hence $x^n \in \mathfrak{p}$. But since $\mathfrak{p}$ is prime, a boring
induction argument gets you that $x \in \mathfrak{p}$. Hence $\sqrt{I} \subset
\mathfrak{p}$ and $\mathfrak{p} \in V(\sqrt{I})$.
\item Let $f \in R \setminus \sqrt{I}$. Then $f^n \notin I$ for all $n$. Hence
$S = \{1, f, f^2, \ldots\}$ is a multiplicative subset, not containing $0$.
Take a
prime ideal $\bar{\mathfrak{p}} \subset S^{-1}R$ containing $S^{-1}I$. Then the
pull-back $\mathfrak{p}$ in $R$ of $\bar{\mathfrak{p}}$ is a prime ideal
containing $I$ that does not intersect $S$. This shows that $\bigcap_{I \subset
\mathfrak p} \mathfrak p \subset \sqrt{I}$. Now if $a \in \sqrt{I}$, then $a^n
\in I$ for some $n$. Hence if $I \subset \mathfrak{p}$, then $a^n \in
\mathfrak{p}$. But since $\mathfrak{p}$ is prime, we have $a \in \mathfrak{p}$.
Thus the equality is shown.
\item $I$ is not the unit ideal if and only if $I$
is contained in some maximal ideal (to
see this, apply (2) to the ring $R/I$) which is therefore prime.
\item If $\mathfrak{p} \in V(I) \cup V(J)$, then $I \subset \mathfrak{p}$ or $J
\subset \mathfrak{p}$ which means that $I \cap J \subset \mathfrak{p}$. Now if
$I \cap J \subset \mathfrak{p}$, then $IJ \subset \mathfrak{p}$ and hence
either $I$ of $J$ is in $\mathfrak{p}$, since $\mathfrak{p}$ is prime.
\item $\mathfrak{p} \in \bigcap_{a \in A} V(I_a) \Leftrightarrow
I_a \subset \mathfrak{p}, \forall a \in A \Leftrightarrow
\mathfrak{p} \in V(\bigcup_{a\in A} I_a)$
\item If $\mathfrak{p}$ is a prime ideal and $f \in R$, then either $f \in
\mathfrak{p}$ or $f \notin \mathfrak{p}$ (strictly) which is what the disjoint
union says.
\item If $a \in R$ is nilpotent, then $a^n = 0$ for some $n$. Hence $a^n \in
\mathfrak{p}$ for any prime ideal. Thus $a \in \mathfrak{p}$ as can be shown by
induction and $D(a) = \emptyset$. Now, as shown in (7), if $a \in R$ is not
nilpotent, then there is a prime ideal that does not contain it.
\item $f \in \mathfrak{p} \Leftrightarrow uf \in \mathfrak{p}$, since $u$ is
invertible.
\item If $\mathfrak{p} \notin V(I)$, then $\exists f \in I \setminus
\mathfrak{p}$. Then $f \notin \mathfrak{p}$ so $\mathfrak{p} \in D(f)$. Also if
$\mathfrak{q} \in D(f)$, then $f \notin \mathfrak{q}$ and thus $I$ is not
contained in $\mathfrak{q}$. Thus $D(f) \cap V(I) = \emptyset$.
\item If $fg \in \mathfrak{p}$, then $f \in \mathfrak{p}$ or $g \in
\mathfrak{p}$. Hence if $f \notin \mathfrak{p}$ and $g \notin \mathfrak{p}$,
then $fg \notin \mathfrak{p}$. Since $\mathfrak{p}$ is an ideal, if $fg \notin
\mathfrak{p}$, then $f \notin \mathfrak{p}$ and $g \notin \mathfrak{p}$.
\item $\mathfrak{p} \in \bigcup_{i \in I} D(f_i) \Leftrightarrow \exists i \in
I, f_i \notin \mathfrak{p} \Leftrightarrow \mathfrak{p} \in \Spec(R)
\setminus V(\{f_i\}_{i \in I})$
\item If $D(f) = \Spec(R)$, then $V(f) = \emptyset$ and
hence $fR = R$, so $f$ is a unit.
\end{enumerate}
\end{proof}

\noindent
The lemma implies that the subsets $V(T)$ from
Definition \ref{definition-spectrum-ring} form the closed
subsets of a topology on $\Spec(R)$. And it also shows that
the sets $D(f)$ are open and form a basis for this
topology.

\begin{definition}
\label{definition-Zariski-topology}
Let $R$ be a ring.
The topology on $\Spec(R)$ whose closed sets are the
sets $V(T)$ is called the {\it Zariski} topology. The open
subsets $D(f)$ are called the {\it standard opens} of $\Spec(R)$.
\end{definition}

\noindent
It should be clear from context whether we consider $\Spec(R)$
just as a set or as a topological space.

\begin{lemma}
\label{lemma-spec-functorial}
\begin{slogan}
Functoriality of the spectrum
\end{slogan}
Suppose that $\varphi : R \to R'$ is a ring homomorphism.
The induced map
$$
\Spec(\varphi) : \Spec(R') \longrightarrow \Spec(R),
\quad
\mathfrak p' \longmapsto \varphi^{-1}(\mathfrak p')
$$
is continuous for the Zariski topologies. In fact, for any
element $f \in R$ we have
$\Spec(\varphi)^{-1}(D(f)) = D(\varphi(f))$.
\end{lemma}

\begin{proof}
It is basic notion (\ref{item-inverse-image-prime}) that
$\mathfrak p := \varphi^{-1}(\mathfrak p')$
is indeed a prime ideal of $R$. The last assertion
of the lemma follows directly from the definitions,
and implies the first.
\end{proof}

\noindent
If $\varphi' : R' \to R''$ is a second ring homomorphism
then the composition
$$
\Spec(R'')
\longrightarrow
\Spec(R')
\longrightarrow
\Spec(R)
$$
equals $\Spec(\varphi' \circ \varphi)$. In other
words, $\Spec$ is a contravariant functor from the
category of rings to the category of topological spaces.

\begin{lemma}
\label{lemma-spec-localization}
Let $R$ be a ring. Let $S \subset R$ be a multiplicative subset.
The map $R \to S^{-1}R$ induces via the functoriality of $\Spec$
a homeomorphism
$$
\Spec(S^{-1}R)
\longrightarrow
\{\mathfrak p \in \Spec(R) \mid S \cap \mathfrak p = \emptyset \}
$$
where the topology on the right hand side is that induced from the
Zariski topology on $\Spec(R)$. The inverse map is given
by $\mathfrak p \mapsto S^{-1}\mathfrak p = \mathfrak p(S^{-1}R)$.
\end{lemma}

\begin{proof}
Denote the right hand side of the arrow of the lemma by $D$.
Choose a prime $\mathfrak p' \subset S^{-1}R$ and let $\mathfrak p$
the inverse image of $\mathfrak p'$ in $R$. Since $\mathfrak p'$
does not contain $1$ we see that $\mathfrak p$ does not contain
any element of $S$. Hence $\mathfrak p \in D$ and we see that
the image is contained in $D$. Let $\mathfrak p \in D$.
By assumption the image $\overline{S}$ does not contain $0$.
By basic notion (\ref{item-localization-zero})
$\overline{S}^{-1}(R/\mathfrak p)$ is not the zero ring.
By basic notion (\ref{item-localize-ideal}) we see
$S^{-1}R / S^{-1}\mathfrak p = \overline{S}^{-1}(R/\mathfrak p)$
is a domain, and hence $S^{-1}\mathfrak p$ is a prime.
The equality of rings also shows that the inverse image of
$S^{-1}\mathfrak p$ in $R$ is equal to $\mathfrak p$,
because $R/\mathfrak p \to \overline{S}^{-1}(R/\mathfrak p)$
is injective by basic notion (\ref{item-localize-nonzerodivisors}).
This proves that the map $\Spec(S^{-1}R) \to \Spec(R)$
is bijective onto $D$ with inverse as given.
It is continuous by Lemma \ref{lemma-spec-functorial}.
Finally, let $D(g) \subset \Spec(S^{-1}R)$ be a standard
open. Write $g = h/s$ for some $h\in R$ and $s\in S$.
Since $g$ and $h/1$ differ by a unit we have $D(g) =
D(h/1)$ in $\Spec(S^{-1}R)$.
Hence by Lemma \ref{lemma-spec-functorial} and the bijectivity
above the image of $D(g) = D(h/1)$ is $D \cap D(h)$.
This proves the map is open as well.
\end{proof}

\begin{lemma}
\label{lemma-standard-open}
Let $R$ be a ring. Let $f \in R$.
The map $R \to R_f$ induces via the functoriality of
$\Spec$ a homeomorphism
$$
\Spec(R_f) \longrightarrow D(f) \subset \Spec(R).
$$
The inverse is given by $\mathfrak p \mapsto \mathfrak p \cdot R_f$.
\end{lemma}

\begin{proof}
This is a special case of Lemma \ref{lemma-spec-localization}.
\end{proof}

\noindent
It is not the case that every ``affine open'' of a
spectrum is a standard open. See
Example \ref{example-affine-open-not-standard}.

\begin{lemma}
\label{lemma-spec-closed}
Let $R$ be a ring. Let $I \subset R$ be an ideal.
The map $R \to R/I$ induces via the functoriality of
$\Spec$ a homeomorphism
$$
\Spec(R/I) \longrightarrow V(I) \subset \Spec(R).
$$
The inverse is given by $\mathfrak p \mapsto \mathfrak p / I$.
\end{lemma}

\begin{proof}
It is immediate that the image is contained in $V(I)$.
On the other hand, if $\mathfrak p \in V(I)$
then $\mathfrak p \supset I$ and we may consider
the ideal $\mathfrak p /I \subset R/I$. Using
basic notion (\ref{item-isomorphism-theorem}) we see that
$(R/I)/(\mathfrak p/I) = R/\mathfrak p$ is a domain
and hence $\mathfrak p/I$ is a prime ideal. From this
it is immediately clear that the image of $D(f + I)$
is $D(f) \cap V(I)$, and hence the map is a homeomorphism.
\end{proof}



\begin{lemma}
\label{lemma-quasi-compact}
\begin{slogan}
The spectrum of a ring is quasi-compact
\end{slogan}
Let $R$ be a ring. The space $\Spec(R)$ is quasi-compact.
\end{lemma}

\begin{proof}
It suffices to prove that any covering of $\Spec(R)$
by standard opens can be refined by a finite covering.
Thus suppose that $\Spec(R) = \cup D(f_i)$
for a set of elements $\{f_i\}_{i\in I}$ of $R$. This means that
$\cap V(f_i) = \emptyset$. According to Lemma
\ref{lemma-Zariski-topology} this means that
$V(\{f_i \}) = \emptyset$. According to the
same lemma this means that the ideal generated
by the $f_i$ is the unit ideal of $R$. This means
that we can write $1$ as a {\it finite} sum:
$1 = \sum_{i \in J} r_i f_i$ with $J \subset I$ finite.
And then it follows that $\Spec(R)
= \cup_{i \in J} D(f_i)$.
\end{proof}

\begin{lemma}
\label{lemma-topology-spec}
Let $R$ be a ring.
The topology on $X = \Spec(R)$ has the following properties:
\begin{enumerate}
\item $X$ is quasi-compact,
\item $X$ has a basis for the topology consisting of quasi-compact opens, and
\item the intersection of any two quasi-compact opens is quasi-compact.
\end{enumerate}
\end{lemma}

\begin{proof}
The spectrum of a ring is quasi-compact, see
Lemma \ref{lemma-quasi-compact}.
It has a basis for the topology consisting of the standard opens
$D(f) = \Spec(R_f)$
(Lemma \ref{lemma-standard-open})
which are quasi-compact by the first remark.
The intersection of two standard opens is quasi-compact
as $D(f) \cap D(g) = D(fg)$. Given any two quasi-compact opens
$U, V \subset X$ we may write $U = D(f_1) \cup \ldots \cup D(f_n)$
and $V = D(g_1) \cup \ldots \cup D(g_m)$. Then
$U \cap V = \bigcup D(f_ig_j)$ which is quasi-compact.
\end{proof}






